\chapter{Methodology and Research Framework}
\label{chapter:methodology}

This project uses the \emph{Development Oriented Triangulation} (DOT) framework\footnote{\url{https://ictresearchmethods.nl/dot-framework/}} to structure our research. The DOT framework breaks things down into three levels: \emph{what} we're researching (domains), \emph{why} we're researching it (balancing fit, expertise, overview, and certainty), and \emph{how} we're researching it (strategies and methods). Below, we'll map out which research strategies and methods we'll use for each part of the project.

\section{Library Research}

Library research explores what's already known and what tools, theories, and guidelines can help us. It mainly looks at the \emph{available work} domain and focuses on building \emph{expertise}.

\subsection{Literature Study}
We'll conduct a systematic literature study covering:
\begin{itemize}
    \item AI/ML methods for time-series anomaly detection (Isolation Forest, LSTM autoencoders, Prophet, ARIMA, and transformer-based approaches)
    \item eBPF internals, the Cilium architecture, and the Hubble observability layer
    \item Zero-trust network models (NIST SP~800-207) and how they apply to Kubernetes
    \item The CNCF landscape for observability and security tooling
\end{itemize}
This method primarily answers sub-questions~1, 2, and~3.

\subsection{Available Product Analysis}
We'll do a comparative analysis of existing tools in two categories:
\begin{itemize}
    \item \textbf{Monitoring stack:} Prometheus, VictoriaMetrics, Grafana, Loki, OpenTelemetry Collector, Jaeger, and commercial alternatives (Datadog, Dynatrace, New Relic)
    \item \textbf{Security and networking:} Cilium, Calico, Flannel, Falco, Tetragon, and Hubble
\end{itemize}
We'll assess each tool based on open-source license, how well it works with Kubernetes, community maturity, feature coverage, and how easy it is to integrate.

\subsection{Best, Good, and Bad Practices}
We'll review these industry standards and guidelines:
\begin{itemize}
    \item CIS Kubernetes Benchmark (v1.9)
    \item CNCF Cloud Native Security Whitepaper
    \item NIST SP~800-207 (Zero Trust Architecture)
    \item OWASP Kubernetes Security Cheat Sheet
\end{itemize}

\subsection{IT Architecture Sketching}
Using ArchiMate~3.2 notation, we'll produce architecture views based on the ArchiMate 3.2 specification, TOGAF ADM documentation, and reference architectures from CNCF projects:
\begin{itemize}
    \item \textbf{Motivation view:} Goals, requirements, and constraints
    \item \textbf{Application layer view:} Monitoring components and their data flows
    \item \textbf{Technology layer view:} Kubernetes objects, Cilium, and the AI inference pipeline
    \item \textbf{Implementation and migration view:} Deployment phases
\end{itemize}
This method uses documented architectural patterns and standards from available work, addressing sub-question~4.

\section{Field Research}

Field research explores the \emph{application domain} and focuses on making sure the solution actually \emph{fits} its operational context.

\subsection{Problem Analysis}
We'll do a structured analysis of the monitoring and security gaps in a representative, freshly deployed Kubernetes cluster (using kubeadm with three worker nodes). The analysis will document:
\begin{itemize}
    \item How many alerts fire per day using only threshold-based rules
    \item Network flows that are invisible to the default networking stack
    \item Incidents that would go undetected without AI-based detection
\end{itemize}
This method directly answers sub-question~1.

\subsection{Domain Modelling}
We'll create an ArchiMate domain model representing the enterprise Kubernetes infrastructure: nodes, namespaces, deployments, services, ingress controllers, storage classes, and external dependencies. This model gives us a shared vocabulary for the architecture design phase.

\section{Lab Research}

Lab research tests whether the solution actually works as intended. It focuses on \emph{certainty}.

\subsection{Prototyping}
We'll use an iterative approach to implement and refine each major component, following official documentation and established practices:
\begin{enumerate}
    \item Cilium installation and basic NetworkPolicy enforcement
    \item Hubble Relay and Hubble UI for flow visualization
    \item Prometheus + Grafana stack with custom Kubernetes dashboards
    \item AI anomaly detection pipeline (training on synthetic normal-traffic baseline, inference via REST API)
    \item Integration layer: Hubble flow exporter $\rightarrow$ OpenTelemetry Collector $\rightarrow$ AI feature store
\end{enumerate}
Key resources include official project documentation (Cilium, Kubernetes, Prometheus), GitHub repositories with example implementations, and community best practices.

\subsection{System Test}
An end-to-end system test will verify that the full monitoring and security stack operates correctly: metrics get scraped, network flows are captured, anomaly detection triggers on injected anomalies, and alerts reach the notification channel.

\subsection{Security Test}
We'll validate Cilium network policies by trying traffic flows that should be blocked (like cross-namespace access without an explicit allow policy, or egress to an unauthorized external host). Tetragon will verify syscall-level enforcement.

\subsection{Non-Functional Test}
We'll measure these non-functional properties:
\begin{itemize}
    \item CPU and memory overhead from Cilium/eBPF compared to a Flannel baseline
    \item Monitoring pipeline latency (time from metric generation to alert firing)
    \item AI inference throughput and model prediction latency under load
\end{itemize}

\section{Showroom Research}

\subsection{Benchmark Test}
We'll benchmark the AI anomaly detection component against a rule-based baseline using precision, recall, F1-score, false-positive rate, and MTTD (mean time to detect). We'll inject anomaly scenarios in a controlled test environment.

\subsection{Guideline Conformity Analysis}
The final implementation will be assessed against the CIS Kubernetes Benchmark and CNCF Security Whitepaper recommendations. We'll report findings as a conformity scorecard in the final research report.

\section{Framework Triangulation}

Table~\ref{tab:triangulation} shows how we're using different research strategies to answer each sub-question.

\begin{table}[H]
\centering
\caption{DOT Framework triangulation: research methods per sub-question.}
\label{tab:triangulation}
\begin{tabular}{@{}p{5.5cm}p{8cm}@{}}
\toprule
\textbf{Sub-question} & \textbf{Primary methods} \\
\midrule
SQ1 --- Limitations of current approach & Problem analysis (Field), Literature study (Library) \\
SQ2 --- AI/ML technique selection       & Literature study (Library), Available product analysis (Library) \\
SQ3 --- Cilium vs.\ traditional networking & Literature study (Library), Benchmark test (Showroom) \\
SQ4 --- Integration architecture        & IT architecture sketching (Library), Prototyping (Lab) \\
SQ5 --- Validation                      & System test (Lab), Security test (Lab), Non-functional test (Lab), Guideline conformity (Showroom) \\
\bottomrule
\end{tabular}
\end{table}

\section{Additional Frameworks}

Beyond the DOT framework, we're also using these methodological frameworks:

\begin{itemize}
    \item \textbf{TOGAF ADM (Architecture Development Method):} The architecture design phase follows the TOGAF ADM cycle, moving from Architecture Vision through Business, Application, and Technology Architecture to the Opportunities and Solutions phase.
    \item \textbf{ArchiMate 3.2:} We're using this as the modeling notation for all architecture views, ensuring interoperability and a standardized vocabulary.
    \item \textbf{GitOps (Flux CD):} All infrastructure-as-code artifacts are managed in a Git repository and automatically reconciled to the cluster, ensuring reproducibility of the proof-of-concept environment.
\end{itemize}
